\chapter{DESARROLLO DEL EMPRENDIMIENTO PRODUCTIVO}
\label{chap:desarrollo_emprendimiento}

En este capítulo se abordan las dimensiones operativas, comerciales y organizacionales que componen la puesta en marcha de la empresa.

\section{Localización del emprendimiento}
\label{sec:localizacion}
Ubicación geográfica, macro y microlocalización de la planta o punto de servicio, argumentando la cercanía a insumos y clientes.

\section{Análisis del mercado}
\label{sec:analisis_mercado}

\subsection{Oferta}
Identificación de los competidores directos e indirectos presentes en el mercado objetivo.

\subsection{Demanda}
Estimación del volumen de consumo y patrón de compra de los clientes potenciales.

\subsection{Público objetivo (cliente y/o usuario)}
Segmentación demográfica, geográfica y socioeconómica del cliente ideal.

\subsection{Entorno y competencia}
Evaluación de las fuerzas del entorno (PESTEL) y análisis de barreras de entrada.

\subsection{Ventaja competitiva del emprendimiento}
Propuesta de valor única que diferencia al producto o servicio de la competencia.

\begin{table}[htbp]
    \centering
    \caption{Matriz sintética de análisis del mercado (Oferta vs. Demanda).}
    \label{tab:estudio_mercado}
    \begin{tabular}{lrr}
        \toprule
        \textbf{Segmento / Zona} & \textbf{Demanda Estimada (Unid/mes)} & \textbf{Oferta Actual (Unid/mes)} \\
        \midrule
        Zona Central            & 5000 & 3200 \\
        Zona Norte              & 3500 & 2100 \\
        Zona Sur                & 4200 & 2800 \\
        \midrule
        \textbf{Total Mercado}  & \textbf{12 700} & \textbf{8100} \\
        \bottomrule
    \end{tabular}
\end{table}


\section{Estrategia de promoción y distribución}
\label{sec:estrategia_promocion}
Canales de distribución, estrategias de precio, marketing digital y publicidad para la captación de clientes.

\section{Estructura organizacional}
\label{sec:estructura_organizacional}
Definición del organigrama, roles, manual de funciones y responsabilidades clave.

\section{Diseño de producto o servicio}
\label{sec:diseno_producto}

\subsection{Características del producto o servicio}
Especificaciones técnicas, empaque, marca y diferenciadores de calidad.

\section{Análisis y descripción del ciclo de producción o de servicio (opcional)}
\label{sec:ciclo_produccion}
Flujograma del proceso productivo o protocolo de atención al cliente desde el ingreso de insumos hasta la entrega final.
